% Transcription — fonds Alexandre Grothendieck, shelfmark 154, batch 7 (pages 121-140).
% Established from the facsimile archives/batches/154/batch-07.pdf, cut from a
% copy of 154.pdf fetched through the project's own relay
% (https://grothendieck-relay.onrender.com/source/154.pdf); PDF page k =
% archivists' page 120+k, no cover sheet. Per the skill transcribe-grothendieck.
% The folder is 175 pages in nine batches, transcribed in parallel; this is
% the seventh. Batch 6 (transcripts/154/batch-06.fr.tex) was read first; its
% notation (and batch 5's) is kept; notation otherwise follows folder 80
% (Systèmes de pseudo-droites).
%
% Pass: Opus 5.5 (claude-opus-5-5), 2026-09-26 — first pass, unchecked against the pages by a human.
%
% Dating and title. \dating is the folder's own date, copied verbatim from
% src/content/catalogue.ts; the folder belongs to the inventory group
% « Espaces stratifiés » (150-151), but since it has a date of its own the
% group's range is not added. Title from the catalogue, trailing full stop
% dropped; « [Système de pseudo-droites] » is bracketed there because the
% archivists supplied it.
%
% Paper. Every leaf is fan-fold computer listing paper. Printed banners
% (« DATASETS PICTURES » tables on 121, 125, 126; JES2 job logs of 19 August
% 1982 on 127, 130, 131, 136, some upside down) are machine dates of the
% paper, not his, and are not transcribed.
%
% Pages skipped: none. Every leaf carries drawings or text of his.
% Pages summarised: none. Pages given only as a figure described in a French
% \note, with their few labels set: 121, 122, 123, 124, 125, 127, 128, 129,
% 133, 138, 140; figures with some formulas or legends: 126, 130, 131, 136,
% 137 (with a sentence), 139.
% Pages transcribed from his hand: 132, 134, 135 (continuous run on the
% « fuseaux » and the « repères » of X~ and the operations sigma_1, sigma,
% sigma_2, sigma_0), 137 (one sentence), 139 (the action of the sigmas on
% a+ a- b+ b-).
%
% His pagination and dates. No page carries a number or a date of his.
% Numbered runs are internal to drawings only: the circuit D_1 ... D_9 = D_1
% on p. 132 (with the reverse circuit), D_1 ... D_20 on p. 136, the vertex
% numberings 1 ... 11 = 1 and 1..4, 1'..4' on p. 131, the circled lines
% 1 ... 11 and letters A ... G on p. 140. p. 131 does not visibly continue
% p. 120; the batch opens on new drawings.
%
% What the batch is about. The unfolded surface X~ of relative positions
% and its cell structure. A « fuseau » (spindle) is the bigon swept by a
% critical relative position with a fixed vertex (p. 132 NB); the circuit
% of stable positions around a spindle depends on (D, s, rho, omega),
% equivalently (D, varpi, omega), varpi an orientation of X at s. Four
% involutions sigma_1, sigma, sigma_2, sigma_0 act on the « repères » of X~
% (pp. 132-135, 138-139): sigma_2 fixes varpi and changes the orientation of
% D; the other is the central symmetry; sigma_1 corresponds to local
% spindles. Dictionary (p. 135): faces of X~ = stable p.r., edges =
% semi-stable p.r., biorientated edges = repères (D, s, omega, varpi),
% vertices = critical p.r.; p. 134 foot: faces <-> spindles, vertices <->
% lines, edges <-> local angle bisectors. p. 137: repères of the surface =
% oriented p.r. of speciality >= 2 with a vertex and an orientation there.
% p. 139: on the model a+ a- b+ b- (1 vertex, 1 face, 2 edges, 4 repères)
% sigma_0 = sigma_2 swaps a+/a-, b+/b-, sigma_1 swaps a+/b+, a-/b-, sigma
% swaps a+/b-, b+/a-; « subdivision barycentrique ». p. 130: counts
% (alpha+gamma+1)/2, (beta+gamma+1)/2 and an inequality. The rest are
% coloured drawings: discs of pseudo-lines swept by fans of positions,
% figure-eight double covers with +/- signs, superposed dual graphs.
%
% Notation fixed once for the batch: \widetilde{X} for his X with tilde (the
% surface of positions, unfolded); \varpi for his barred omega (orientation
% of X at a vertex), \omega for the orientation of D, \Omega for the
% orientations of a spindle; \sigma, \sigma_0, \sigma_1, \sigma_2 for his
% barred sigmas; « p.r. » for position relative; the exponent on X~ read
% « rep » (repères) is doubtful and noted where it occurs.
%
% Readings to check first: the right column of p. 132 (the last four lines
% are very fast); p. 134 left column (« épinglage », the word before
% « fournis ») and the exponent of X~ on pp. 134-135; p. 130 first line; the
% top list on p. 139; the sequence of vertices on p. 136.

\documentclass[11pt,a4paper]{article}
% Le préambule commun à toutes les transcriptions du fonds Grothendieck.
%
% Il tient en une idée : **l'incertitude se compose, elle ne se résout pas.**
% Une transcription qui lisse un mot douteux a détruit la seule chose pour
% laquelle on la faisait — un lecteur doit pouvoir savoir, à chaque endroit,
% ce qui était sur le feuillet et ce qui a été ajouté. Les six macros qui
% suivent sont donc l'appareil critique, et non de la décoration.
%
% Les mêmes commandes sont reconnues par `scripts/render.mjs`, qui produit la
% vue de lecture du site. Une macro ajoutée ici doit l'être aussi là-bas,
% faute de quoi le rendu échouera bruyamment — ce qui est voulu : un rendu
% silencieusement faux est pire qu'un rendu absent.

\ProvidesPackage{grothendieck}[2026/08/08 Transcriptions du fonds Grothendieck]

\usepackage{amsmath,amssymb,amsthm}
\usepackage{mathrsfs}
\usepackage{stmaryrd}
% Les diagrammes commutatifs sont partout dans ce fonds : c'est la langue
% même de la géométrie algébrique de Grothendieck, et une transcription qui
% les rendrait en prose ne transcrirait rien.
\usepackage{tikz-cd}
% Français par défaut — c'est la langue des feuillets — mais l'anglais est
% chargé aussi : la traduction et le résumé sont des documents anglais, et la
% typographie française (espaces avant les deux-points) y serait une faute.
% Ces documents basculent par \selectlanguage{english} en tête de corps.
\usepackage[english,french]{babel}
\usepackage{fontspec}
\usepackage{geometry}
\geometry{a4paper,margin=25mm,marginparwidth=28mm}
\usepackage{marginnote}
\usepackage{xcolor}
% ulem, not soul. `soul`'s \st and \ul are built on a character-by-character
% scan that XeLaTeX's font machinery breaks: the document fails with an
% undefined control sequence rather than degrading. ulem does the same two jobs
% and is Unicode-engine safe. `normalem` stops it hijacking \emph, which the
% transcriptions use for Grothendieck's own emphasis.
\usepackage[normalem]{ulem}
\usepackage{hyperref}
\hypersetup{colorlinks=true,linkcolor=black,urlcolor=black,pdfborder={0 0 0}}

% --- Métadonnées du lot -----------------------------------------------------
% Reprises telles quelles par le script de rendu, qui les affiche en tête de
% la vue de lecture. Elles fixent ce que le fichier prétend couvrir : sans
% elles, une transcription flotte sans point d'attache dans un fonds de seize
% mille feuillets.

\newcommand{\folder}[1]{\gdef\grothFolder{#1}}
\newcommand{\batch}[1]{\gdef\grothBatch{#1}}
\newcommand{\leaves}[2]{\gdef\grothFirst{#1}\gdef\grothLast{#2}}
\newcommand{\foldertitle}[1]{\gdef\grothFoldertitle{#1}}
\gdef\grothFolder{?}\gdef\grothBatch{?}\gdef\grothFirst{?}\gdef\grothLast{?}\gdef\grothFoldertitle{}

% --- Le résumé --------------------------------------------------------------
%
% Il ouvre la lecture modernisée, sous le titre « Résumé », et il est composé
% en petit. Ce n'est pas un ornement : c'est le seul passage du document qui
% ne soit pas des mathématiques, et un lecteur doit pouvoir le reconnaître
% avant de l'avoir lu — puis le sauter s'il connaît déjà le sujet, ou s'y
% arrêter s'il ne le connaît pas. Le filet qui le suit marque la reprise du
% corps.

\newenvironment{resume}
  {\small\setlength{\parskip}{0.45em}}
  {\par\medskip\hrule\medskip}

% --- Mots-clés ---------------------------------------------------------------
%
% En anglais, en clôture du résumé de la lecture modernisée : le vocabulaire
% moderne sous lequel chercher ce que les feuillets construisent. Le manifeste
% du site (`npm run manifest`) les extrait pour étiqueter la cote entière —
% cette ligne est la source unique des tags, il n'y a pas de fichier à côté.
\newcommand{\keywords}[1]{\par\smallskip\noindent
  {\footnotesize\textsc{Keywords} --- \textit{#1}}\par}

% --- L'appareil critique ----------------------------------------------------

% Le numéro de feuillet, en marge. C'est l'ancre qui permet de régler un
% désaccord sur une formule en regardant *un* feuillet plutôt qu'une chemise
% de six cents. Le site s'en sert aussi pour tourner les pages du fac-similé
% quand on fait défiler la transcription.
\newcommand{\leaf}[1]{%
  \marginnote{\footnotesize\color{gray!70}\textsf{#1}}%
  \label{leaf:#1}\ignorespaces}

% Illisible : on ne devine pas. Un mot inventé qui se lit comme les autres est
% le pire résultat possible.
\newcommand{\ill}{\textcolor{red!60!black}{[\,\ldots\,]}}

% Lecture douteuse : le mot est donné, et signalé comme incertain.
\newcommand{\uncertain}[1]{\uline{#1}}

% Ajout de l'éditeur — une lettre manquante, un mot sous-entendu.
\newcommand{\add}[1]{\textcolor{blue!50!black}{[#1]}}

% Biffé par Grothendieck lui-même. Ce qu'il a rayé fait partie du document :
% une rature dit souvent où la pensée a changé de direction.
\newcommand{\struck}[1]{\sout{#1}}

% Note du transcripteur, sur l'état matériel du feuillet.
\newcommand{\note}[1]{\par\smallskip{\small\color{gray!60!black}\textit{[#1]}}\par\smallskip}

% Note marginale de Grothendieck, distincte du corps du texte.
\newcommand{\marginal}[1]{\par\smallskip\noindent\hspace{1em}%
  {\small\color{gray!60!black}#1}\par\smallskip}

% --- Titre ------------------------------------------------------------------

\AtBeginDocument{%
  \begin{center}
    {\large\bfseries Fonds Alexandre Grothendieck --- cote n\textsuperscript{o}~\grothFolder}\\[2pt]
    {\itshape\grothFoldertitle}\\[4pt]
    {\small feuillets \grothFirst--\grothLast\ (lot \grothBatch)}\\[2pt]
    {\footnotesize\color{gray!60!black}Transcription établie d'après le fac-similé numérisé
     par l'Université de Montpellier.\\
     Les lectures incertaines sont soulignées, les illisibles notées [\,\ldots\,],
     les ajouts entre crochets.}
  \end{center}
  \vspace{1em}}

\endinput


\folder{154}
\batch{7}
\pages{121}{140}
\dating{1983-1984}
\watermark{Édition de démonstration}
\foldertitle{[Système de pseudo-droites] : notes manuscrites (1983-1984, s.d.)}

\begin{document}

\page{121}
\note{page de dessins en couleur, sans texte. En haut, quatre petits rectangles bordés de vert ou de brun, traversés de droites rouges qui se croisent. Un disque bordé de vert, rempli de pseudo-droites rouges, traversé par deux droites vert clair fléchées marquées $D_1$ et $D_\nu$ ; un disque rouge rempli de courbes rouges et vertes qui vont d'un point du bord au point opposé. Au crayon, trois droites $D_2$, $D_0$, $D_1$ qui passent à travers deux croisements orange. Au milieu, une droite vert clair qui contourne des croisements orange, les sommets marqués $s_p$, $s_0$, $s_1$, $s_2$, et les étiquettes $D_p$, $D_1$, $D_\nu$. En bas, un hexagone rouge au pourtour épais, rempli de tirets orange, aux sommets duquel se croisent des droites rouges (deux points marqués) ; une grille de droites rouges et orange ; une bande de droites vert clair coupée de croisements orange, avec un chemin en tirets}

\page{122}
\note{à gauche, un disque bordé de jaune, deux points marqués sur le bord, contenant une courbe rouge qui se recoupe plusieurs fois ; à côté, le mot « \emph{Non} » (souligné). À droite, des croquis au crayon : un cercle traversé d'une courbe, des droites qui se croisent en étoile sur une droite, des arcs}

\page{123}
\note{dessins en couleur, sans texte : une pseudo-droite rouge qui coupe une droite jaune, une région hachurée au crayon ; un disque bordé de rouge, deux points jaunes opposés sur le bord, traversé de pseudo-droites rouges et d'un faisceau de fines courbes au crayon qui joignent les deux points (un balayage) ; un demi-disque rouge dont la base jaune porte un segment noir ; un cercle au crayon ; un grand disque aux pseudo-droites rouges et noires très enchevêtrées, balayé de même par un faisceau de courbes au crayon entre deux points du bord ; enfin, un croisement orange posé sur une droite verte}

\page{124}
\note{dessins en couleur, sans texte : des droites rouges qui se croisent sur une droite vert clair fléchée, avec des arcs noirs qui contournent un croisement (un balayage à travers un triangle), les arcs successifs numérotés $1$ à $5$ ; au crayon, une suite de flèches numérotées $1$ à $5$ ; plusieurs croisements rouges sur des droites vertes et brunes ; un disque bordé de vert traversé de droites rouges, deux arcs rouges épais sur le bord ; en bas, deux petits disques verts, l'un traversé d'un fuseau d'arcs rouges joignant deux points opposés du bord, l'autre d'un diamètre rouge ; des disques au crayon}

\page{125}
\note{dessins en couleur, sans texte : un éventail d'arêtes orange (pleines et en tirets) issu d'un sommet, dans un pentagone bleu et vert ; une étoile bleue dont le centre est entouré de pétales bruns ; un réseau cellulaire vert (sommets d'ordre $4$) et bleu, avec en son centre la même fleur de pétales et une bande bleue hachurée ; un cercle orange traversé d'une étoile de rayons bleus ; un petit quadrillage bleu et vert ; un triangle curviligne noir contenant une étoile verte, avec un arc hachuré}

\page{126}
\note{en haut, une bande : une droite vert clair entre deux droites brunes, coupée de croisements orange ; les segments de la droite supérieure sont marqués $p'$, $q'$, $s'$, $u'$, $v'$, $w'$, $y'$, $z'$, ceux de l'inférieure $p$, $q$, $r$, $t$, $u$, $w$, $x$, $z$, avec de petites marques $x$, $1$, $2$, $3$ et des barres noires sur certains croisements. À droite, un cercle au crayon traversé de droites orange, aux extrémités marquées de barres bleues ; des morceaux de graphe vert (en H, en Y) ; un heptagone bleu avec des cordes orange et une étoile verte au centre. En bas, des étoiles au crayon avec des flèches, une fleur de pétales hachurés et un grand « $2$ » ; un petit peigne orange}

\[
\boxed{\sigma_1,\ \sigma,\ \sigma_2,\ \sigma_0}
\]
\note{encadré au milieu du bas de la page}

\page{127}
\note{dessins en couleur, sans texte : une colonne de trois disques tangents (bleus et au crayon) traversés par une même pseudo-droite orange ; un disque rempli de pseudo-droites orange, entouré d'un cercle vert en tirets, tangent à un cercle vert vide ; un disque aux pseudo-droites orange dont le bord porte de petites encoches bleues en demi-cercle ; un disque vert et bleu rempli de pseudo-droites orange ; un disque au crayon ; un polygone au crayon avec une étoile de flèches orange et de rayons bleus}

\page{128}
\note{dessins en couleur, sans texte : trois graphes superposés, bleu, vert et orange, autour d'un sommet bleu central, avec six petits carrés bruns aux sommets de l'hexagone orange ; une étoile bleue à huit branches ; un sommet bleu entouré d'un cercle orange hachuré ; un graphe vert en tirets ; un hexagone vert portant des triangles orange à ses sommets, traversé de rayons bleus}

\page{129}
\note{dessins en couleur, sans texte sinon des signes : plusieurs paires de disques tangents (figures en huit) traversés de pseudo-droites rouges, des triangles coloriés en jaune ; dans la paire centrale, des points marqués $+$ dans le disque du haut et $-$ dans celui du bas, et des flèches doubles près du point de contact ; un réseau vert traversé de droites orange ; des disques au crayon, l'un au bord jaune ondulé}

\page{130}
\note{à gauche, dessins en couleur : deux disques tangents traversés de pseudo-droites rouges, un arc du bord hachuré en noir ; un disque vert traversé de pseudo-droites rouges, le bord marqué d'accolades $\alpha$, $\gamma$, $\beta$ ; un disque vert aux pseudo-droites rouges ; à droite, deux cercles, rouge et vert, qui se coupent en deux points}

$\alpha + \gamma + 1$ \ill{}
\[
\frac{\alpha + \gamma + 1}{2} \qquad \frac{\beta + \gamma + 1}{2}
\]
\[
\frac{\alpha + \beta + 2\gamma + 2}{2} \leqslant \alpha + \beta + \gamma + 1 \quad \ill{}
\]
\struck{$\alpha + \beta + \gamma + 1$} \qquad \struck{$3 - \dfrac{\alpha + \beta}{2} \leqslant$}

$\dfrac{\alpha + \beta}{2}$ \ill{} \ill{} \ill{} $(\geqslant 1)$
\note{la lecture du premier terme de la ligne biffée est incertaine ; le « $\geqslant$ » final est surchargé}

\page{131}
\note{dessins en couleur : plusieurs disques verts ou orange traversés de pseudo-droites orange ; dans deux d'entre eux, un arc vert clair (une pseudo-droite en position relative) et les points d'intersection numérotés sur le bord $1, 2, \ldots, 10$, avec « $11 = 1$ », et dans un troisième, bordé de violet, les points $1, 2, 3, 4$ et $1', 2', 3', 4'$ ; un demi-disque brun au bord orné de petites flèches ; des cercles noirs fléchés ; un rectangle au crayon aux coins numérotés $1, \ldots, 4$ et $1', \ldots, 4'$ ; un cercle au crayon dont le bord porte $\xi_0, \xi_1, \ldots, \xi_{n-2}, \xi_{n-1}$ et $\xi'_1, \xi'_2, \ldots, \xi'_{n-2}$}

\[
2n \big/ 2(n+1) \qquad 2 \geqslant 2n
\]
\note{le second terme est surchargé ; au-dessous, une droite vert clair traversée d'une flèche verticale}

\[
\begin{array}{cccccccc}
1 & 2 & 3 & 4 & 1' & 2' & 3' & 4' \\
\xi_0 & \xi_1 & \xi_2 & \xi_3 = \xi_{n-1} & & & &
\end{array}
\]

\[
\widetilde{\Delta}_{\tilde{s}} \xrightarrow[\sim]{\varphi_0} \widetilde{\Delta}'_{\tilde{s}_0}
\qquad
\widetilde{\Delta}_{\tilde{s}} \xrightarrow{\varphi_1} \widetilde{\Delta}^{\ast}_{\tilde{s}_1}
\]
\note{les deux flèches partent du même $\widetilde{\Delta}_{\tilde{s}}$, l'une horizontale, l'autre en diagonale ; les indices $\tilde{s}_0$, $\tilde{s}_1$ sont d'une lecture incertaine}

\[
\alpha,\ u^{n}\xi_0,\ u^{n+1}\xi_0,\ \ldots,\ u^{2(n-1)-1}\xi_0
\qquad
\alpha,\ u^{n}\xi_0,\ u^{n+1}\xi_0,\ \ldots,\ u^{2n-2}\xi_0
\]
\[
\xi_0\ \xi_1 \cdots \xi_{n-2}\ \xi_{n-1} = \xi'_0
\qquad
\xi_0\ \xi_1 \cdots \xi_{n-1}\ \xi_n = \xi'_0
\]
\note{les deux groupes de formules sont écrits de part et d'autre d'un cercle vert ; sous le premier, « $u^{n-1}$ » ; les exposants sont d'une lecture incertaine}

\[
\begin{array}{cccc}
1 & 2 & 3 & 4 \\
1' & 2' & 3' & 4'
\end{array}
\qquad
0\ 2'\ 3'\ 4'\ 0'\ 1\ 2\ 3\ 4\ 1'
\]
\note{dans le premier tableau, $2, 3, 4$ et $2', 3', 4'$ sont encadrés et une flèche ramène de $4$ à $1'$ ; dans la suite de droite, des arcs relient $2'$ à $2$, $3'$ à $3$, $4'$ à $4$ et $0$ à $1'$}

$\operatorname{Circ}\bigl(\widetilde{I \setminus \lbrace i_0 \rbrace}\bigr) \leftarrow \tilde{\imath}_0$
\note{au crayon ; en bas à droite, un cercle noir traversé d'un diamètre, les extrémités marquées $\alpha$ et $\alpha'$, une flèche de rotation}

\page{132}
\note{en haut à gauche, un dessin en couleur : un fuseau (bigone) bordé d'un arc vert foncé et d'un arc vert clair, entre deux croisements orange, traversé de pseudo-droites orange et de cinq courbes brunes numérotées $1$ à $5$ aux deux bouts, qui le balaient ; au milieu, un second dessin du même fuseau, les courbes numérotées $1$ à $8$, les croisements marqués « $1 = 5$ », « $2$ », « $3$ », « $4$ », « $5 = 8$ » (en orange), « $5/6$ » sur l'arc supérieur et « $1/2$ » sur l'inférieur}

\marginal{\emph{fuseau} avec : $a$ des deux rives (= côtés du bigone) du fuseau ; $\varpi$ des deux extrémités ; $\Omega$ des deux orientations du fuseau ; $\boxed{\Omega \simeq a \wedge \varpi}$}
\note{en haut à droite, les quatre lignes réunies par une accolade ; « (= côtés du bigone) » est écrit au-dessus de la ligne, « du fuseau » ajouté après « orientations » ; la lettre $\varpi$ est écrite sur une autre, biffée}

$\Omega' \simeq \Omega \wedge \varpi \simeq a$ \qquad $\Omega'$ une des orientations de $X$ en le \add{sommet du fuseau}. Si on a choisi un $D \in a$, donc $a \simeq \lbrace \pm 1 \rbrace$, \struck{\ill{}} $\Omega \simeq \varpi \simeq \omega$ (une des orientations de $D$)
\note{« sommet du fuseau » est écrit au-dessus de la ligne ; la lettre lue $\varpi$ dans « $\Omega \wedge \varpi$ » peut aussi être un $\sigma$}

\emph{NB} « Fuseau » $\longleftrightarrow$ p.r.\ critique, avec sommet fixé dessus (« él.\ \ill{} de contact »)

On trouve donc le circuit $D_1, D_2, \ldots, D_8, D_9 = D_1$. Comment trouvons-nous le circuit inverse $D_2, D_1 = D_9, D_8, \ldots, D_3, D_2$ ? C'est en changeant \emph{à la fois} la rive $\rho$, et l'orientation $\omega$.

Ainsi, l'opération $\sigma_2$ sur les repères de $\widetilde{X}$ correspond à l'opération qui \struck{\ill{}} \emph{fixe} $\varpi$ (orientation en $s$) et qui change \add{la} coorientation de $D$, \struck{i.e.\ des} \add{\uncertain{ou}} \ill{} du fuseau)
\note{au-dessous, un croquis au crayon : un cercle partagé en secteurs marqués $D_1, \ldots, D_8$, un axe fléché passant par le centre}

Ce circuit de positions relatives stables dépend de la position semi-stable de départ $D$ \add{avec le sommet $s \in S_D$} (vert clair), du choix de sa générisation de départ, i.e.\ de la rive correspondante $\rho$ de $D$ en $s$, et d'une \struck{\ill{}} orientation $\omega$ de $D$ (qui indique l'ordre de succession \struck{des repères} \ill{} des triangles), ce qui (moyennant le choix de $\rho$) équivaut à la donnée d'une orientation $\varpi$ de $X$ en $s$.
\note{colonne de droite ; à côté, un petit dessin : un croisement orange sur la droite vert clair $D$, orientée par $\omega$, le sommet $s$ et, au-dessous, la rive brune $\rho$}

Donner $(D, \varpi, \omega)$ revient \uncertain{donc} \uncertain{à} \ill{} que de donner le fuseau qu'il définit, avec un \struck{\ill{}} \ill{} du fuseau $\Longleftrightarrow$ une orientation $\varpi$ \struck{de} $X$ en son sommet, et une \uncertain{somme} (\uncertain{déployée}) du fuseau i.e.\ une de ses \ill{}
\note{les quatre dernières lignes, serrées au bas de la page, sont d'une lecture très incertaine}

\page{133}
\note{croquis au crayon, sans texte : des arrangements de trois ou quatre droites, et un cercle coupé de quatre droites, dont les régions portent le nombre de leurs côtés ($3$, $4$, $5$) ; un petit cercle traversé de courbes (au crayon, et un autre en vert)}

\page{134}
\struck{\ill{}} en somme, le choix de $(D, \varpi, \omega)$ revient : \uncertain{elle} au fuseau correspondant, avec un \uncertain{épinglage} de ce fuseau i.e.\ des bigônes qu'il définit.

De ce point de vue, $\sigma_2$ est l'opération qui fixe \struck{\ill{}} l'élément de structure \add{\uncertain{de ce type}} de dim $1$ (\ill{} \ill{}), et fixe celui de dim $0$. Pour $\sigma_0$, serait-ce l'inverse ?

\emph{Non}, cette opération inverse, \struck{i.e.} en tant qu'opération $\sigma$ dans les repères de $\widetilde{X}$, est celle qui remplace un repère par son symétrique p.r.\ au \add{sommet} \struck{\ill{}} origine, qui commute bien à $\sigma_2$,
\[
\begin{array}{ccccc}
\sigma & - & \sigma_2 & - & \sigma_0 \\
\backslash & & & & \\
\sigma_1 & & & &
\end{array}
\]
\note{petit graphe à gauche de la ligne suivante}
On a les quatre opérations \struck{fuseaux locaux} fuseaux globaux $\leftarrow$ des $\sigma_1, \sigma, \sigma_2, \sigma_0$ sur les repères, dont \struck{deux} covariantes commutant. On vient d'interpréter $\sigma, \sigma_2$ comme opérations sur les fuseaux, il resterait à interpréter $\sigma_1$ et $\sigma_0$.
\note{« fuseaux locaux » est biffé et souligné d'une accolade ; « fuseaux globaux » est écrit à la suite, une accolade sous « globaux »}

Pour $\sigma_1$, on trouve que les classes de $\widetilde{X}^{\mathrm{rep}}/\sigma_1$, i.e.\ les secteurs angulaires de $\widetilde{X}$, correspondent aux \ill{} fournis par une p.r.\ stable (\ill{})
\note{l'exposant de $\widetilde{X}$, lu « rep » (repères), est incertain ; de même page 135}

\note{colonne de droite : deux dessins en couleur. Un fuseau bordé de tirets vert clair et vert foncé, entre deux croisements orange, balayé d'une courbe au crayon qui revient sur elle-même, une flèche vers la droite ; un fuseau semblable, prolongé des deux côtés par une droite noire, les sommets marqués $a$ (avec l'orientation $\omega_a$) et $b$ (avec $\omega_b$), de petites flèches de rotation}

et deux triangles formés « \struck{anti}covariants », donnant naissance à \emph{deux} p.r.\ sous-stables donnant \uncertain{avec} une orientation \struck{\ill{}} \add{et} orientation \struck{\ill{}} \add{de} $X$ en le sommet $a$ (resp.\ $b$), l'une de l'autre étant induite par une orientation du bigone vert : \struck{C'est} \add{les} fuseaux locaux, i.e.\ des $\sigma_1$ et la symétrie centrale, \struck{celle de $\sigma_2$}. Il faudrait comprendre aussi les classes de $\widetilde{X}^{\mathrm{rep}}/(\sigma_1, \sigma)$, qui sont fournis \uncertain{dans} \ill{} $\lbrace D_{1/2}, D_{2/3}, D_{5/6}, D_{6/7}$
\note{« anti » est biffé dans « anticovariants », le mot souligné ; l'accolade ouverte n'est pas refermée}

\note{en bas à droite, un dessin : deux croisements orange sur une droite noire, entourés de tirets bruns et d'une courbe au crayon qui les enjambe ; légendes « arête de $D$ » et « p.r.\ critique $D$ »}

\uncertain{faces} $\leftrightarrow$ \uncertain{fuseaux}

sommets $\leftrightarrow$ droites

arêtes $\leftrightarrow$ bisect.\ ang.\ locaux
\note{au bas de la page à gauche, en petits caractères}

\page{135}
Sur $\widetilde{X}^{\mathrm{rep}}$, il y a les opérations $\sigma_1, \sigma, \sigma_2, \sigma_0$ dont deux covariantes commutent.
\note{en haut à droite, deux petits croquis au crayon d'un triangle posé sur une droite, le second marqué $\sigma$, et une flèche verticale}

\emph{faces de $\widetilde{X}$} : p.r.\ stables ; \emph{faces orientées} : p.r.\ \uncertain{orientées} stables

\emph{arêtes de $\widetilde{X}$} : p.r.\ semi-stables ; \struck{arêtes biorientées}

\emph{arêtes biorientées i.e.\ repères} \emph{ou repères} : \struck{\ill{}} \add{repère} $(D, s, \omega, \varpi)$ où : $D$ p.r.\ critique ; $s \in S_D = D \cap S$ ; $\omega$ orientation de $D$, $\varpi$ or.\ de $X$ en $s$
\note{les trois conditions sont en colonne, derrière un trait vertical ; « repère » est écrit au-dessus du mot biffé ; le $\varpi$ est ajouté au-dessus de la parenthèse}

i.e.\ \emph{Donner repère p.r.} : $(D, s_D)$, $+$

\emph{sommets} : p.r.\ \emph{critiques}

\emph{sommets orientés} : p.r.\ \emph{orientées} critiques

repères $\to$ faces orientées : $(D, s, \omega) \longmapsto$ la position relative stable obtenue en \struck{\ill{}}
\note{un $\varpi$ est écrit en surcharge devant la parenthèse ; la phrase s'arrête là ; le reste de la page est vide}

\page{136}
\note{dessins en couleur : en haut, un grand fuseau entre deux croisements orange, bordé d'un arc orange, traversé de droites vert clair, la droite vert foncé et deux chemins en tirets bruns qui passent par les sommets ; à gauche, un petit repère fléché ($\omega$, $\varpi$). À droite, un cercle au crayon traversé de deux paires de droites orange, une droite vert clair et une vert foncé fléchée, les sommets $s_4$, $t_3$ (chiffres entourés), les étiquettes $D_1 = D_{19}$, $D_{18}$, $D_4$, et au croisement de droite « $5 = 19$ ! ». Au milieu, un croisement orange sur une droite verte avec des flèches $\varpi$, $\omega$ ; à gauche en bas, un croisement enjambé d'une courbe au crayon, avec la légende « \uncertain{non sans triangle} ». En bas, un grand fuseau entre deux croisements orange, bordé d'un arc orange, les sommets marqués $s_1$, $s_2$, $s_3$, $s_4$, $s_5$, $s_6 = s_1$, $t_2$, $t_3$, avec des numéros bleus et noirs $1$ à $19$ autour des sommets ; à droite, une ellipse au crayon traversée de croisements orange et de la droite vert foncé}

$P_f = {}$
\[
s_1\ s_2\ s_3\ s_4\ s_5\ s_6 = s_1\ s_2\ t_2\ s_3,\ s_4,\ t_3,\ s_5,\ s_6 = s_1,\ s_2,\ t_2,\ t_3,\ s_4,\ t_3,\ s_5 \quad s_6 = s_1
\]
\[
D_1 \mid D_2\ D_3\ D_4\ D_5\ D_6 \mid D_7\ D_8\ D_9\ D_{10}\ D_{11}\ D_{12}\ D_{13}
\]
\[
\mid D_{14}\ D_{15}\ D_{16}\ D_{17}\ D_{18}\ D_{19}\ D_{20} \overset{?}{=} D_6 \mid
\]
\note{sur la page, les deux suites sont écrites l'une au-dessus de l'autre, sur une seule ligne chacune ; les barres verticales sont des traits de couleur (vert clair, vert foncé, vert, noir) ; une flèche sous $D_4\,D_5$ ; sous $D_{18}$ et $D_{19}$, « $D_4 \to D_5$ » relié par des signes $=$ ; dans la suite des sommets, le troisième $t$ de la dernière série est surchargé, et l'alignement des sommets sur les $D_i$ est approximatif}

\page{137}
\note{dessins en couleur : un grand demi-disque orange rempli de pseudo-droites orange, avec, le long de sa base, un fuseau (tirets vert foncé, droite vert clair, tirets bruns) entre deux sommets marqués $1$, $2$ et des flèches de rotation ; en haut à droite, une bande étroite (vert foncé et vert clair) coupée de croisements orange, les arcs numérotés $1$ à $5$, et aux bouts « $2 : 5$ » ; une étoile au crayon avec deux sommets entourés (vert et vert foncé) et des tirets numérotés $1$, $2$ ; des croquis au crayon ; un croisement orange sur un chemin en tirets bruns ; en bas à droite, un long fuseau entre deux croisements orange, bordé de vert foncé et de vert clair, traversé de croisements orange de plus en plus serrés, un chemin en tirets bruns par les sommets, avec la légende « fuseau de croisement »}

$\sigma_2 = {}$\struck{é}changement de \struck{\ill{}} \add{l'orientation} de $D$ \add{et}

$\sigma_0 = {}$?
\note{au crayon, en bas à gauche ; « l'orientation » est écrit au-dessus du mot biffé et souligné, « et » au-dessous de la ligne}

les repères de la surface correspondent aux p.r.\ \add{$D$ orientées} qui \add{sont de spécialité} \struck{\ill{}} $\geqslant 2$, \emph{munies} de plus d'un sommet\supplied{$s$} et d'une orientation en ce sommet i.e.\ d'une rive du complémentaire $D \setminus \lbrace s \rbrace$
\note{au crayon ; « $D$ orientées » et « sont de spécialité » sont écrits au-dessus de la ligne, la seconde insertion au-dessus d'un mot biffé ; la dernière lecture, « $D \setminus \lbrace s \rbrace$ », est incertaine}

\page{138}
\note{en haut à gauche, un graphe en couleur (flèches orange et vert clair) : quatre sommets $\sigma_1$, $\sigma_2$, $\sigma$, $\sigma_0$ et, au milieu des arêtes, les sommets $(\sigma_1, \sigma_2)$, $(\sigma_1, \sigma)$, $(\sigma, \sigma_2)$, $(\sigma_2, \sigma_0)$, $(\sigma, \sigma_0)$ et, en haut, $(\sigma_1, \sigma_0)$, relié par un grand arc vert ; chaque sommet-couple reçoit les flèches de ses deux termes et une flèche venue de l'extérieur, marquée du produit correspondant ($\sigma_1\sigma_2$, $\sigma\sigma_1$, $\sigma_0\sigma_2$, $\sigma\sigma_2$ en tirets, $\sigma\sigma_0$, $\sigma_1\sigma_0$). En haut à gauche, la suite « $\sigma_1\ \sigma\ \sigma_2\ \sigma_0$ ». Le reste de la page est vide}

\page{139}
\note{en haut à gauche, un petit dessin : un segment de $-1$ à $+1$ passant par $0$, les deux moitiés marquées $a^{+}$, $b^{+}$ au-dessus et $a^{-}$, $b^{-}$ au-dessous ; au-dessous, le même dessin, biffé, et deux flèches}

\uncertain{$1$ sommet}

$1$ face

$2$ arêtes

\struck{$2$} $4$ repères, correspondant aux \uncertain{rives} des arêtes $a$, $b$
\note{le premier chiffre de la première ligne est surchargé}

$\sigma_0$, \struck{\ill{}} \quad $\sigma_0 = \sigma_2$ échangent $a^{+}, a^{-}$ et $b^{+}, b^{-}$

$\sigma_1$ échange \struck{\ill{}} $a^{+}$ et $b^{+}$, $a^{-}$ et $b^{-}$

$\sigma$ échange $a^{+}$ et $b^{-}$, $b^{+}$ et $a^{-}$
\note{à gauche, « $\sigma_1, \sigma, \sigma_2, \sigma_0$ », les deux premiers et les deux derniers réunis par une accolade, la seconde marquée d'une lettre lue $\sigma'$}

subdivision barycentrique

$\sigma_0 = \sigma_2$ soit $\sigma'$, commute à $\sigma$

$\sigma_1 = \sigma\sigma'$ \quad \struck{$= \sigma_1\sigma'$}

$\sigma_1$
\[
\underbrace{\sigma\sigma'}_{\sigma_1},\ \underbrace{\sigma}_{\sigma},\ \underbrace{\sigma'}_{\sigma_2},\ \underbrace{\sigma'}_{\sigma_0}
\qquad\qquad
\underbrace{\sigma\sigma'}_{\sigma'_0},\ \underbrace{\sigma'}_{\sigma'_1},\ \underbrace{\sigma}_{\sigma'_2}
\]
\note{sous les deux derniers termes de la première liste, une accolade commune}

\note{dessins au crayon et à l'encre brune : une sphère coupée par un équateur en tirets, les points $-1$, $+1$, $0$, $\infty$ et un méridien épais ; au-dessous, une sphère au crayon traversée de méridiens, le point $0$ en haut, avec une légende illisible ; au milieu, un ovale brun (moitié en tirets) aux points $0$, $\infty$, $-1$, $+1$ ; en haut à droite, deux dômes au crayon aux méridiens marqués, avec les points $-1$, $0$, $+1$}

\page{140}
\note{grand dessin au crayon et à l'encre, sans texte : un faisceau de pseudo-droites étiquetées de lettres $A, \ldots, G$ et de numéros entourés $1, 1', 2, \ldots, 11$, dont les points d'intersection sont nommés $a_0, a_1, \ldots, a_{10}$, $b_1, b_2, \ldots, b_{10}$ (avec « $b_{-1} = b_{10}$ » en haut), $c_1, c_{11}$, $d_1, d_{11}$, $e_1, e_{10}, e_{11}$, $f_1, f_{11}$, $g_1, g_{11}$ ; des bandes hachurées et en tirets suivent les pseudo-droites, et trois triangles sont noircis. À droite, un quadrillage de droites $A, \ldots, G$ (verticales) et $1, \ldots, 5$ (horizontales), avec des diagonales tracées sur les colonnes $D$ à $G$ ; au-dessus, une double ligne fléchée numérotée $1$, $2$ ; et deux petits croquis de droites qui se coupent}

\end{document}
